\section{Laws and moments of the critical field}\label{sec:fields}

Section~\ref{sec:correlations} gives the signed correlations at separated
points and a bound valid up to the two-point diagonal.  We now pass from
these estimates to the field limit.  The main issue is to remove
collisions of arbitrarily many lattice sites using only the two-point
bound.  A Lee--Yang estimate for absolute moments supplies this step.
We also construct the continuum field and its mixed moments, so that
the full Dirichlet variance remains visible in the identification.

The argument develops the proofs of \cite[Lemmas~5.10--5.11 and
Theorem~5.12]{B}.  We give the transfer argument here with the exact
normalization of this paper: the common scalar multiplying
$\exp\{G_n(x,x)/(2K_*)\}e^{i\sigma\theta_x}$ is
$a_n^{-1}\exp\{-G_n(0,0)/(2K_*)\}$, which may depend on $n$.
The statement of the cited transfer theorem instead uses a fixed
scalar.  Throughout this section,
\[
 S_n^\sigma(f)=n^{-2}\sum_{x\in D_n^\circ}f(x)W_n^\sigma(x),
 \qquad \sigma\in\{-1,1\},
\]
with $W_n^\sigma$ defined in Section~\ref{sec:correlations}; thus
$S_n^+=S_n$ and
$S_n^-(f)=\overline{S_n^+(\overline f)}$.
The same sign convention will be used for continuum fields.  In
particular the minus sign conjugates the field, and conjugating an
entire evaluation also conjugates its test function.

\subsection{Absolute moments from Lee--Yang}

The following finite-volume estimate applies to arbitrary nonnegative
weights.  Their size, spatial variation, and dependence on the mesh
will therefore place no restriction on its use for $S_n$.

\begin{lemma}[Lee--Yang absolute moment bound]\label{fields:lee-yang}
Consider a finite ferromagnetic XY model with nonnegative edge
couplings and all boundary angles fixed to zero.  For deterministic
weights $\lambda_x\geq0$, set
\[
 W=\sum_x\lambda_x e^{i\theta_x},\qquad v=\E|W|^2.
\]
There is an absolute constant $C$ such that
\begin{equation}\label{fields:ly-bound}
 \|W\|_{L^p}\leq C\sqrt p\,v^{1/2},\qquad p\geq2.
\end{equation}
For every finite list of vertices, including repetitions, and every
choice of signs $\sigma_i\in\{-1,1\}$, one also has
\begin{equation}\label{fields:positive-moments}
 \E\exp\left\{i\sum_i\sigma_i\theta_{x_i}\right\}\geq0.
\end{equation}
\end{lemma}

\begin{proof}
Let $U$ be independent and uniform on the circle, and rotate every
spin, including each boundary spin, by $U$.  Subtracting the common
boundary angle recovers the original fixed-boundary law.  Put
$X=\operatorname{Re}(e^{iU}W)$.  The finite-graph XY Lee--Yang
Theorem, with the common uniformly distributed boundary convention,
says that the entire function $\E e^{zX}$ has only imaginary zeros
\cite[Theorem~1 and Remark~5]{NewmanWu2019}.
One may obtain this boundary convention on simple graphs by adjoining
a vertex of observable weight zero, joining it ferromagnetically to
every boundary vertex, and taking all these added couplings to infinity.
The limiting spins on the boundary equal the new spin, whose angle is
uniform by rotation invariance.  The bounded observable gives locally
uniform convergence of the Laplace transforms; the zero-free
half-planes persist under this limit.

The random variable $X$ is symmetric and bounded.  Hence its
Lee--Yang product representation
\cite[Proposition~13]{NewmanWu2019} is
\[
 \E e^{zX}
   =e^{Bz^2}\prod_j(1+z^2/y_j^2),\qquad
 B\geq0,\quad \sum_jy_j^{-2}<\infty,\quad
 B+\sum_jy_j^{-2}=\tfrac12\E X^2.
\]
The zeros are listed as pairs $\pm i y_j$, with $y_j>0$ and
multiplicity.  Since uniform rotation gives $\E X^2=v/2$, the
inequality $1+u\leq e^u$ yields
\begin{equation}\label{fields:ly-laplace}
 \E e^{tX}\leq e^{t^2v/4},\qquad t\in\R.
\end{equation}
The variable $Y=\operatorname{Im}(e^{iU}W)$ has the same law as $X$.
If $v>0$, Chernoff's Inequality applied to both signs of $X$ and $Y$
therefore gives
\[
 \Pbb\{|W|>r\}
 \leq \Pbb\{|X|>r/\sqrt2\}+\Pbb\{|Y|>r/\sqrt2\}
 \leq4e^{-r^2/(2v)},\qquad r>0.
\]
Integrating this bound against $p r^{p-1}\,dr$ gives
$\E|W|^p\leq2p(2v)^{p/2}\Gamma(p/2)$, which implies
Equation~\eqref{fields:ly-bound}.  When $v=0$, the assertion is
immediate.  The auxiliary rotation has left $|W|$ unchanged.

For Equation~\eqref{fields:positive-moments}, orient the edges
arbitrarily and expand each interaction in its absolutely convergent
Fourier series.  If an edge has coupling $b\geq0$, its coefficients are
\[
 I_k(b)=\sum_{r=0}^\infty
       \frac{(b/2)^{2r+|k|}}{r!(r+|k|)!}\geq0,
 \qquad k\in\Z.
\]
At $b=0$ these are interpreted as $I_0(0)=1$ and $I_k(0)=0$
for $k\ne0$.  Integrating each interior angle imposes the divergence
constraint prescribed by the inserted charges on the integer edge
currents.  Boundary angles are zero, so they contribute a factor one
and impose no divergence constraint.  Each surviving summand is
nonnegative, and the partition function in the denominator is
positive.  Repeated vertices simply add their integer charges.
\end{proof}

\subsection{The continuum circle field}

We next identify the continuum moments which arise from
Equation~\eqref{eq:source-17}.  We state the construction for the
whole range in which the collision bound below is integrable.
The circle-averaging argument follows \cite[Lemma~5.11]{B};
its use of covariance positivity is related to the Onsager estimate
in \cite[Proposition~3.9]{JSW2020}.  We include the bound for unequal
radii, which is needed to compare two regularizations.

\begin{lemma}[Circle Wick field and its mixed moments]
\label{fields:circle-chaos}
Let $K>1/(4\pi)$ and let $\Phi_D$ be the zero-Dirichlet Gaussian free
field with covariance $G_D=(-\Delta_D)^{-1}$.  Where the circle of
radius $\varepsilon$ about $x$ lies in $D$, define
\[
 V_\varepsilon(x)=
 \exp\left\{\frac{i\Phi_{D,\varepsilon}(x)}{\sqrt K}
        +\frac{\E\Phi_{D,\varepsilon}(x)^2}{2K}\right\}.
\]
There exists a complex random distribution
$V_{K,D}\in H^{-2}_{\mathrm{loc}}(D)$ such that, for every
$f\in C_c^\infty(D;\C)$,
$V_\varepsilon(f)\to V_{K,D}(f)$ in every finite $L^p$.
For every smooth compactly supported cutoff $\chi$,
$\chi V_\varepsilon\to\chi V_{K,D}$ in
$L^2(\Omega;H^{-2})$.
For $m\geq1$, arbitrary signs $\sigma_i\in\{-1,1\}$, and
$f_i\in C_c^\infty(D;\C)$,
\begin{equation}\label{fields:continuum-moments}
 \E\prod_{i=1}^m V_{K,D}^{\sigma_i}(f_i)
 =\int_{D^m}\prod_{i=1}^m f_i(x_i)
    \exp\left\{-\frac1K\sum_{i<j}\sigma_i\sigma_j
                         G_D(x_i,x_j)\right\}
                  \,dx_1\cdots dx_m.
\end{equation}
The integral is absolutely convergent.  In particular,
$\E V_{K,D}(f)=\int_D f(x)\,dx$.
\end{lemma}

\begin{proof}
Fix a compact set $E\Subset D$ containing the supports of the tests,
and choose $0<\rho<1$ so that its closed $2\rho$-neighborhood lies
in $D$.  On that neighborhood write
\[
 G_D(u,v)=\frac1{2\pi}\log\frac1{|u-v|}+R_D(u,v),
 \qquad q=\frac1{2\pi K}<2.
\]
The remainder $R_D$ is smooth and bounded, including across the
diagonal.  Let $C_{a,b}(x,y)$ be the covariance of the circle
averages of $\Phi_D/\sqrt K$ with centers $x,y\in E$ and radii
$a,b\leq\rho$.  We first establish, uniformly in both radii,
\begin{equation}\label{fields:circle-covariance}
 C_{a,b}(x,y)
   =q\log\frac1{\max\{|x-y|,a,b\}}+O_{E,K}(1).
\end{equation}
To see this, put $r=|x-y|$ and $M=\max\{r,a,b\}$.  The double
circle average of the logarithm of the distance equals
\[
 L_{a,b}(r)=\frac1{2\pi}\int_0^{2\pi}
                 \log\max\{a,|r+be^{it}|\}\,dt.
\]
Here we have used the elementary circle-mean identity
$(2\pi)^{-1}\int_0^{2\pi}\log|z-ae^{it}|\,dt
=\log\max\{|z|,a\}$.
The integrand is at most $\log(2M)$.  If $a=M$, it is at least
$\log M$.  Otherwise, its integral is bounded below by the circle
mean of $\log|r+be^{it}|$, which is $\log\max\{r,b\}=\log M$.
Thus $\log M\leq L_{a,b}(r)\leq\log M+\log2$.
Together with the bounded remainder this proves
Equation~\eqref{fields:circle-covariance}, even for coincident centers.
The finite logarithmic integrals also justify the covariances of
the circle averages.

For arbitrary radii $\varepsilon_i\leq\rho$, Gaussian integration
and cancellation of the full diagonal variances give
\begin{equation}\label{fields:regularized-kernel}
 \E\prod_{i=1}^m V_{\varepsilon_i}^{\sigma_i}(x_i)
   =\exp\left\{-\sum_{i<j}\sigma_i\sigma_j
                  C_{\varepsilon_i,\varepsilon_j}(x_i,x_j)\right\}.
\end{equation}
At the pointwise level $V_\varepsilon^-$ denotes
$\overline{V_\varepsilon}$.  To dominate this expression for all
signs and all collision patterns, suppose that $m\geq2$ and the
centers are distinct, and set
\[
 d_i=\min\{\rho,\min_{j\ne i}|x_i-x_j|\},
 \qquad r_i=\max\{\varepsilon_i,d_i/4\}.
\]
All circles with radii $r_i$ remain in $D$.  For $i\ne j$,
\[
 \max\{|x_i-x_j|,r_i,r_j\}
   =\max\{|x_i-x_j|,\varepsilon_i,\varepsilon_j\}.
\]
Consequently enlarging the radii changes each off-diagonal
covariance by a bounded amount.  If $C^r$ is the covariance matrix
for the enlarged circles, its positive semidefiniteness gives
\[
 -\sum_{i<j}\sigma_i\sigma_j C^r_{ij}
   =\frac12\sum_i C^r_{ii}
       -\frac12\sum_{i,j}\sigma_i\sigma_j C^r_{ij}
   \leq\frac12\sum_i C^r_{ii}.
\]
Since $C^r_{ii}\leq q\log(4/d_i)+C_{E,K}$, we obtain
\begin{equation}\label{fields:continuum-domination}
 0<\E\prod_i V_{\varepsilon_i}^{\sigma_i}(x_i)
     \leq C_{E,K,m}\prod_i d_i^{-q/2}.
\end{equation}
This bound is independent of every radius, including their
relative sizes.  For $m=1$ the expectation is exactly one.

We verify explicitly that the right side is integrable on $E^m$.
Put $\alpha=q/2<1$ and use
\[
 d_i^{-\alpha}
    \leq\rho^{-\alpha}+\sum_{j\ne i}|x_i-x_j|^{-\alpha}.
\]
Expanding the product produces finitely many directed graphs on
$\{1,\ldots,m\}$, with at most one outgoing edge from each vertex
and no self-edge.  An edge $i\to j$ contributes
$|x_i-x_j|^{-\alpha}$.  Each connected component is a tree,
or consists of one directed cycle with trees attached.
Vertices with no incoming edge can be integrated first: each has
at most one remaining incident edge, and
\begin{equation}\label{fields:polar-integral}
 \sup_{y\in E}\int_E |x-y|^{-s}\,dx<\infty,\qquad 0\leq s<2.
\end{equation}
This follows by polar integration over a fixed disk containing
$E-E$.  It removes all trees at a uniform cost.  A cycle on two
vertices contributes $|x-y|^{-2\alpha}$ and is integrable because
$2\alpha=q<2$.  For a cycle on at least three vertices, integrate
one vertex using Cauchy--Schwarz:
\[
 \int_E |x-y|^{-\alpha}|x-z|^{-\alpha}\,dx
 \leq
 \left(\int_E|x-y|^{-2\alpha}\,dx\right)^{1/2}
 \left(\int_E|x-z|^{-2\alpha}\,dx\right)^{1/2}
 \leq C_{E,\alpha}.
\]
The remaining graph is a path, which can be integrated from its
ends.  Isolated vertices contribute $|E|$.  Hence every term of
the expansion is integrable, including near simultaneous
collisions of any number of centers.

At distinct centers,
$C_{\varepsilon_i,\varepsilon_j}(x_i,x_j)\to G_D(x_i,x_j)/K$
as all radii tend to zero, with no restriction on their relative
rates.  Equations~\eqref{fields:regularized-kernel}
and~\eqref{fields:continuum-domination} now give convergence of
all integrated mixed kernels by dominated convergence.
In particular, the second mixed kernels with radii
$\varepsilon,\delta$ show that $V_\varepsilon(f)$ is Cauchy
in $L^2$.  Taking equal numbers of positive and negative signs
gives uniform bounds for every even absolute moment.
For any finite $p$, choose an even order greater than $p$.
These bounds and $L^2$ convergence give convergence in $L^p$
by uniform integrability.  H\"older's Inequality then passes
each mixed product to the limit and proves
Equation~\eqref{fields:continuum-moments}.

It remains to realize these limits as one random distribution.
Choose a flat torus whose fundamental square contains $\overline D$
in its interior, with normalized Fourier basis
$e_\nu$, $\nu\in\Z^2$.  For a fixed
$\chi\in C_c^\infty(D)$, Equation~\eqref{fields:continuum-domination}
with two factors implies
\[
 \sup_{\nu,\ 0<\varepsilon<\varepsilon_\chi}
       \E|V_\varepsilon(\chi e_\nu)|^2<\infty
\]
for sufficiently small $\varepsilon_\chi>0$.
Each Fourier coefficient is Cauchy in $L^2$, and
$\sum_{\nu\in\Z^2}(1+|\nu|^2)^{-2}<\infty$.
Dominated convergence in this series proves that
$\chi V_\varepsilon$ is Cauchy in $L^2(\Omega;H^{-2})$ on
the torus.  Choose a countable exhaustion by smooth cutoffs
which equal one near the supports of their predecessors.
Their limits are compatible, since multiplication by a smooth
function is continuous on $H^{-2}$.  They define an almost surely
$H^{-2}_{\mathrm{loc}}$ distribution $V_{K,D}$, and testing
recovers the scalar limits already constructed.  Localization
identifies these torus Sobolev spaces with the Euclidean local
spaces on $D$.  This proves all the asserted convergence.
In particular the normalization uses the complete Dirichlet
variance, including its finite spatial part.
\end{proof}

At $K=K_*=2/\pi$ the parameter $q$ in the proof is $1/4$,
and Lemma~\ref{fields:circle-chaos} constructs precisely the
field $V_*$ in Theorem~\ref{thm:main}.  We next prove that its
mixed moments are the limits of the lattice moments.

\subsection{Uniform bounds and removal of collisions}

\begin{lemma}[Bounds for smeared lattice fields]
\label{fields:moment-bounds}
For every compact $E\Subset D$, every bounded complex function $f$
supported in $E$, and every $p\geq2$,
\begin{equation}\label{fields:test-bounds}
 \E_n^0|S_n^\sigma(f)|^2\leq C_E\|f\|_\infty^2,
 \qquad
 \|S_n^\sigma(f)\|_{L^p}\leq C_E\sqrt p\,\|f\|_\infty,
\end{equation}
uniformly in $n$ and $\sigma\in\{-1,1\}$.
If $B\subset E$ is a square of side $d\leq1$, then
\begin{equation}\label{fields:box-bounds}
 \E_n^0|S_n^\sigma(1_B)|^2\leq C_E(d+n^{-1})^3,
 \qquad
 \|S_n^\sigma(1_B)\|_{L^p}
       \leq C_E\sqrt p\,(d+n^{-1})^{3/2}.
\end{equation}
\end{lemma}

\begin{proof}
Write $h=n^{-1}$.  For a fixed $x\in B\cap D_n^\circ$,
the shell $jh\leq|x-y|<(j+1)h$ contains at most $C(j+1)$
lattice points.  Equation~\eqref{eq:source-20} thus implies
\[
 h^2\sum_{y\in B\cap D_n^\circ}(|x-y|+h)^{-1}
   \leq Ch\sum_{0\leq j\leq C(d/h+1)}1
   \leq C(d+h).
\]
There are at most $C(d/h+1)^2$ choices of $x$, so a second
summation, including its factor $h^2$, proves the first bound
in Equation~\eqref{fields:box-bounds}.  Summation over a fixed
compact set gives a uniform bound for its indicator.  Positivity
of the opposite-sign two-point expectation then gives, for
arbitrary complex $f$ supported there,
\[
 \E_n^0|S_n^+(f)|^2
 \leq h^4\sum_{x,y\in E\cap D_n^\circ}
       |f(x)|\,|f(y)|\,
            \E_n^0[W_n^+(x)W_n^-(y)]
 \leq C_E\|f\|_\infty^2.
\]

For a nonnegative weight $w$, the coefficients of the spins in
$S_n^+(w)$ are
\[
 \lambda_x=n^{-2}w(x)a_n^{-1}
       \exp\left\{\frac{G_n(x,x)-G_n(0,0)}{2K_*}\right\}\geq0.
\]
Lemma~\ref{fields:lee-yang} applies to these exact coefficients
for each $n$.  It proves the second box bound from the first.
For a complex weight, decompose
\[
 f=(\operatorname{Re}f)_+-(\operatorname{Re}f)_-
       +i(\operatorname{Im}f)_+-i(\operatorname{Im}f)_-.
\]
Each nonnegative summand is bounded by $\|f\|_\infty$ and
supported in $E$.  Minkowski's Inequality and the Lee--Yang
bound prove the second estimate in
Equation~\eqref{fields:test-bounds}.  Conjugation proves the
same bounds for the minus field.
\end{proof}

\begin{proposition}[Convergence of every mixed moment]
\label{fields:moments}
For every $m\geq1$, signs $\sigma_i\in\{-1,1\}$, and
$f_i\in C_c^\infty(D;\C)$,
\begin{equation}\label{fields:lattice-moment-limit}
 \lim_{n\to\infty}\E_n^0\prod_{i=1}^m S_n^{\sigma_i}(f_i)
       =\E\prod_{i=1}^m V_*^{\sigma_i}(f_i).
\end{equation}
The right side is the integral in
Equation~\eqref{fields:continuum-moments} with $K=K_*$.
\end{proposition}

\begin{proof}
Let $E\Subset D$ contain all supports.  The case $m=1$
follows from Equation~\eqref{eq:source-17}, uniformly on $E$,
and Riemann summation.  Assume $m\geq2$ and expand the moment
as a sum over $(x_1,\ldots,x_m)\in(D_n^\circ)^m$.
All signed correlations in this sum are nonnegative by
Equation~\eqref{fields:positive-moments} and the positivity
of the deterministic normalizers.

Fix two indices $i\ne j$ and a small $d>0$.  Enlarging the squares
of a grid of side $d$ by a fixed factor gives a family
$\mathcal B_d$ of at most $C_Ed^{-2}$ squares, each of side
at most $C d$, such that every pair $x_i,x_j\in E$ with
$|x_i-x_j|<d$ lies in one of these squares.  All these squares
lie in a fixed larger compact set $E'\Subset D$ when $d$
is small.  Positivity allows us to bound the absolute value
of the restricted moment sum by
\[
 \left(\prod_{\ell=1}^m\|f_\ell\|_\infty\right)
 \sum_{B\in\mathcal B_d}
 \E_n^0\left[
     S_n^{\sigma_i}(1_B)S_n^{\sigma_j}(1_B)
       \prod_{\ell\ne i,j}S_n^{\sigma_\ell}(1_{E'})
       \right].
\]
Each expectation is nonnegative.  Its absolute value is at
most the expectation of the product of absolute values.
H\"older's Inequality with $m$ equal exponents and
Lemmas~\ref{fields:lee-yang} and~\ref{fields:moment-bounds} bound it by
\[
 C_{E',m}
 \|S_n^+(1_B)\|_{L^2}^2
 \leq C_{E',m}(d+n^{-1})^3.
\]
Indeed the two box factors each contribute their
$L^2$ norm times a constant depending only on $m$, and
every other factor has a uniformly bounded $L^m$ norm.
Summing over all pairs proves
\begin{equation}\label{fields:collision-bound}
 \limsup_{n\to\infty}
 n^{-2m}\!\!\sum_{\substack{x_1,\ldots,x_m\in D_n^\circ\\
                    \min_{i<j}|x_i-x_j|<d}}
 \left(\prod_i|f_i(x_i)|\right)
          \E_n^0\prod_i W_n^{\sigma_i}(x_i)
 \leq C_{E,m,f_1,\ldots,f_m}d.
\end{equation}
In particular repeated lattice sites, of every multiplicity
and every choice of signs, are included in this estimate.

Choose a continuous cutoff on configurations which vanishes
when some pair is at distance at most $d$ and equals one
when every pair is at distance at least $2d$.
On its support, the uniform separated-point convergence in
Equation~\eqref{eq:source-17} turns the lattice moment sum
into the Riemann sum of the continuum kernel.
The cutoff makes the product with the test functions continuous
across the excluded diagonals.  For fixed $d$ these sums
therefore converge to their continuum integrals.
Equation~\eqref{fields:collision-bound} bounds the omitted
lattice contribution by a quantity tending to zero with $d$,
after taking the upper limit in $n$.
The omitted continuum contribution tends to zero by the
integrable domination in
Equation~\eqref{fields:continuum-domination}.
Letting $d\downarrow0$ proves
Equation~\eqref{fields:lattice-moment-limit}.
\end{proof}

\subsection{Identification of laws and local tightness}

The moment bounds also supply the determinacy needed to identify
joint complex laws.  This step uses the real and imaginary parts
together, so it retains the nonneutral moments.

\begin{proposition}[Joint convergence of test-function laws]
\label{fields:law}
For every $k\geq1$ and $f_1,\ldots,f_k\in C_c^\infty(D;\C)$,
\[
 (S_n(f_1),\ldots,S_n(f_k))
    \ \Longrightarrow\
 (V_*(f_1),\ldots,V_*(f_k)).
\]
Every mixed moment of this vector and its complex conjugate
converges.
\end{proposition}

\begin{proof}
Fix a real linear combination $Z_n$ of the real and imaginary
parts of the $k$ lattice evaluations, and let $Z$ be the
same linear combination of the continuum evaluations.
Equation~\eqref{fields:test-bounds} gives
$\|Z_n\|_{L^p}\leq C\sqrt p$ for all $p\geq2$, with $C$
depending on the tests and the chosen linear combination.
Consequently, for every $t\geq0$,
\begin{equation}\label{fields:exponential-moments}
 \sup_n\E_n^0 e^{t|Z_n|}<\infty.
\end{equation}
Indeed the exponential series is dominated, apart from its
first two terms, by
$\sum_{p\geq2}(Ct\sqrt p)^p/p!<\infty$.
Proposition~\ref{fields:moments} gives convergence of the even
moments of $Z_n$ to those of $Z$.  Hence $Z$ satisfies the
same bound at even orders, and at all remaining orders by
H\"older's Inequality.  It too has every exponential moment.
Absolute convergence of the exponential series thus determines
its characteristic function from its moments.

The finite-dimensional lattice vectors are tight by their
second-moment bounds.  Along a weakly convergent subsequence,
the higher uniform moment bounds make every polynomial in
their real and imaginary coordinates uniformly integrable.
The limit therefore has the moments given in
Proposition~\ref{fields:moments}.  Each of its real projections
has the same even-moment bounds as $Z$, hence is exponentially
integrable and has the same characteristic function as $Z$.
The Cram\'er--Wold Criterion identifies the vector law.
Every convergent subsequence has this law, which proves
convergence along the full sequence.  For the stated mixed
moments, use the sign convention
$\overline{S_n(f)}=S_n^-(\overline f)$ in
Proposition~\ref{fields:moments}.
\end{proof}

\begin{proposition}[Tightness in the local Sobolev topology]
\label{fields:tightness}
The laws of $S_n$ are tight in $H^{-3}_{\mathrm{loc}}(D)$.
They converge there, along the full sequence, to the law of $V_*$.
\end{proposition}

\begin{proof}
Use the larger flat torus from the continuum construction and
its normalized Fourier basis $e_\nu$, $\nu\in\Z^2$.
Choose real smooth cutoffs $\chi_j$ with compact support in $D$,
equal to one on an increasing exhaustion of $D$, and with
$\chi_{j+1}=1$ near $\supp\chi_j$.
Since the Fourier functions have uniformly bounded suprema,
Equation~\eqref{fields:test-bounds} gives
\[
 \sup_{n,\nu}\E_n^0|S_n(\chi_j e_\nu)|^2\leq C_j.
\]
Fourier summation yields
\begin{equation}\label{fields:sobolev-bound}
 \sup_n\E_n^0\|\chi_j S_n\|_{H^{-2}}^2
       \leq C'_j\sum_{\nu\in\Z^2}(1+|\nu|^2)^{-2}
       <\infty.
\end{equation}
The embedding $H^{-2}\hookrightarrow H^{-3}$ on this torus
is compact.  Indeed finite Fourier truncations have finite rank,
and on their complements the ratio of the squared Sobolev
weights is $(1+|\nu|^2)^{-1}\to0$.

Given $\eta>0$, choose $R_j<\infty$ so that the sum of the
Markov bounds in Equation~\eqref{fields:sobolev-bound},
with thresholds $R_j^2$, is less than $\eta$.
With probability at least $1-\eta$, all inequalities
$\|\chi_jS_n\|_{H^{-2}}\leq R_j$ hold simultaneously.
The set of distributions satisfying these inequalities has
relatively compact closure in $H^{-3}_{\mathrm{loc}}(D)$:
compact embedding gives a convergent subsequence for each
localized field, and a diagonal choice gives convergence for
all $j$.  The identities
$\chi_j(\chi_{j+1}S_n)=\chi_jS_n$ pass to the limits because
smooth multiplication is continuous on $H^{-3}$.  Thus the
limits define one distribution on $D$.  This proves tightness.

For completeness, the topology just used is the Euclidean
local topology in the statement of Theorem~\ref{thm:main}.
On distributions supported in any fixed compact subset of $D$,
the torus Sobolev norms and the corresponding Euclidean
Sobolev norms are equivalent by localization inside its
fundamental square.  Moreover every compactly supported
cutoff is supported where some $\chi_j$ equals one.
The countable family of these localized norms therefore
generates exactly $H^{-3}_{\mathrm{loc}}(D)$.

Every subsequential limit in this space has the test-function
laws in Proposition~\ref{fields:law}, since evaluation on a
smooth compactly supported test is continuous.
These laws determine the distribution-valued law.  For
example, the countably many tests $\chi_j e_\nu$ determine
every localized Fourier coefficient; their real and imaginary
parts generate the Borel sigma-field of the compatible
subspace of the countable product of the separable spaces
$H^{-3}$.  Lemma~\ref{fields:circle-chaos} places $V_*$
in this space, so all subsequential limits have its law.
Tightness and this uniqueness prove full-sequence convergence.
\end{proof}

Propositions~\ref{fields:moments}, \ref{fields:law},
and~\ref{fields:tightness} complete the proof of
Theorem~\ref{thm:main}.  Every lattice field in this argument
uses the same exact center magnetization $a_n$ and the same
Green-function factor fixed in its definition.
